\title{Controlling Text-to-Image Diffusion by Orthogonal Finetuning}

\begin{abstract}
Recent advances in large-scale text-to-image diffusion models have yielded strikingly realistic image synthesis based on textual descriptions. However, steering or conditioning these robust models to accomplish diverse downstream objectives remains a significant challenge. In response, we propose Orthogonal Finetuning (OFT), a theoretically motivated approach for refining text-to-image diffusion models for task-specific adaptation. OFT, in contrast to prior strategies, is designed to rigorously maintain hyperspherical energy—a property encapsulating the relationships between neuron activations constrained on a unit hypersphere. Our analysis indicates that retaining this energy is fundamental to conserving the semantic image generation proficiency of text-to-image diffusion architectures. To further bolster the robustness of finetuning, we introduce Constrained Orthogonal Finetuning (COFT), which integrates a radius constraint onto the hypersphere to further regularize the adaptation process. In our study, we focus on two salient finetuning scenarios: subject-driven generation, where the aim is to synthesize individualized images from limited exemplars and accompanying text, and controllable generation, which extends the model’s capability to interpret external control cues. Through comprehensive experiments, we demonstrate that our OFT framework not only surpasses baseline methods in both quality of image generation and training efficiency but also achieves superior convergence behavior.
\end{abstract}

\section{Introduction}

State-of-the-art text-to-image diffusion models~\cite{saharia2022photorealistic,ramesh2022hierarchical,rombach2022high} have substantially advanced the field of photorealistic image synthesis conditioned on text prompts. Nevertheless, textual instructions alone often lack precision, making it difficult to exert granular and accurate control over the resulting outputs. This paper focuses on addressing two principal types of text-to-image generative tasks:

\begin{itemize}[leftmargin=*,nosep]

    \item \textbf{Subject-driven generation}~\cite{ruiz2023dreambooth}: Presented with only a handful of sample images depicting a particular subject, the objective is to render new depictions of that subject situated in various novel contexts using textual input as additional guidance.
    \item \textbf{Controllable generation}~\cite{zhang2023adding,mou2023t2i}: Provided with supplementary control information (such as canny edge maps or semantic segmentation layouts), the model’s task is to synthesize images that conform to both this control input and the corresponding textual prompt.
\end{itemize}

The central challenge in both tasks is determining how to finetune text-to-image diffusion models in a way that retains the strong generative capabilities learned during pretraining. To characterize a desirable finetuning approach, we identify two key objectives: (1) \emph{training efficiency}, meaning a reduction in the number of trainable parameters and required training epochs, and (2) \emph{generalizability preservation}, which involves maintaining both high fidelity and diversity in the generated outputs. Commonly, finetuning is performed either by making incremental adjustments to the existing neuron weights with a low learning rate (\emph{e.g.}, \cite{ruiz2023dreambooth}) or by incorporating a compact module with separately parameterized weights (\emph{e.g.}, \cite{hulora2022,zhang2023adding}). Despite the straightforward nature of these techniques, neither ensures that the generative strengths achieved during pretraining are preserved. Furthermore, there remains a gap in rigorous theoretical guidance for crafting robust finetuning methodologies or for selecting key hyperparameters, such as the number of epochs for training. A significant hurdle lies in the absence of a concrete metric that quantifies the extent to which the pretrained generative competence is retained. Most current finetuning approaches are based on the implicit assumption that minimizing the Euclidean distance between the parameters of the finetuned and pretrained models leads to better retention of pretrained abilities, which in turn motivates the use of extremely small learning rates. However, although this belief is sometimes valid, we contend that Euclidean parameter distance alone does not sufficiently reflect semantic preservation. Thus, leveraging a more structurally meaningful metric to evaluate the divergence between finetuned and pretrained models would likely enhance both the retention of pretrained capabilities and the overall stability of finetuning.

Building on empirical findings that hyperspherical similarity effectively captures semantic relationships~\cite{liu2017deep,liu2018decoupled,chen2020angular}, our approach utilizes hyperspherical energy~\cite{liu2018learning} to describe the pairwise interactions among neurons. Specifically, hyperspherical energy, defined as the aggregate of pairwise hyperspherical similarities (\emph{e.g.}, cosine similarity) for neurons within a given layer, serves as a measure of neuron distribution uniformity on the unit hypersphere~\cite{liu2021learning}. We conjecture that an effective finetuning process should aim to keep the hyperspherical energy as consistent as possible with the pretrained configuration. An initial attempt might be to incorporate a regularization term that maintains hyperspherical energy during finetuning; however, this strategy lacks a formal guarantee of minimizing hyperspherical energy deviation. To address this, we exploit the invariance property of hyperspherical energy—namely, pairwise hyperspherical similarity remains unchanged under shared orthogonal transformations of all neurons. Drawing on this property, we introduce Orthogonal Finetuning~(OFT), which enables the adaptation of large text-to-image diffusion models for downstream tasks without altering their hyperspherical energy. The principal concept is to learn an orthogonal transformation, applied uniformly to each layer, thereby maintaining the relative angular structure of neurons. This can be interpreted as recalibrating the canonical basis of each neuron layer. Furthermore, since retaining a small Euclidean distance between the parameters of the finetuned and pretrained models is known to help preserve pretraining performance, we extend OFT to a variant called Constrained Orthogonal Finetuning~(COFT). This variant restricts the finetuned parameters to reside within a hypersphere of fixed radius centered at the pretrained neuron vectors.

The rationale behind employing orthogonal transformations for neuron finetuning is partly inspired by the concept of the 2D Fourier transform, which allows an image to be expressed in terms of its magnitude and phase spectra. Notably, the phase spectrum—encoding the angular relationship between the input and the basis—retains most of the semantic content. Empirically, combining the phase spectrum from an image with a random magnitude spectrum enables the reconstruction of the original image’s semantic structure. This observation indicates that manipulating neuron directions fundamentally alters the semantics of generated images—a core aim in both subject-driven and controllable image synthesis. However, unrestrained alteration of neuron directions risks compromising the pretrained model’s generative quality. To address this, we introduce a constraint: maintaining the pairwise angles between neurons, guided by the hypothesis that these angular relationships encapsulate key neural network knowledge. Building on this insight, we propose the use of a shared orthogonal transformation for neurons within each layer, thereby preserving the hyperspherical energy.

Additionally, our approach takes cues from orthogonal over-parameterized training~\cite{liu2021orthogonal}, where classification neural networks are trained from scratch via orthogonal transformation of an initially random network. Random initialization results in low expected hyperspherical energy, and the strategy in \cite{liu2021orthogonal} is to keep this energy minimal throughout training, which correlates with improved classification generalization~\cite{liu2018learning,lin2020regularizing}. Their results demonstrate that orthogonal transformations afford sufficient flexibility for training generalizable classifiers. In contrast, our focus is on the finetuning of text-to-image diffusion models, aiming to enhance both control and generative performance in downstream tasks. There are two primary distinctions between OFT and the method in \cite{liu2021orthogonal}. First, whereas \cite{liu2021orthogonal} explicitly targets the reduction of hyperspherical energy, OFT seeks to retain the hyperspherical energy of the pretrained model to protect the existing internal structure during finetuning—minimizing this energy for diffusion models would risk erasing essential semantic features. Second, OFT introduces mechanisms to limit deviation from the pretrained model, resulting in a constrained variant, whereas \cite{liu2021orthogonal} operates without such restrictions. Striking an optimal balance between adaptability and stability is fundamental in the finetuning context, and we posit that the OFT framework achieves this balance effectively. Our principal contributions are summarized as follows:

\begin{itemize}[leftmargin=*,nosep]

    \item We introduce Orthogonal Finetuning, a new approach to finetuning text-to-image diffusion models to enhance controllability. Additionally, we present a constrained version of this method, which restricts the angular deviation from the pretrained weights in order to promote further stability.
    \item In contrast with prevailing finetuning strategies, OFT preserves the hyperspherical energy of the model parameters during the finetuning process, a property we demonstrate empirically to be vital for maintaining the generative semantics established by the pretrained model.
    \item We deploy OFT in two distinct scenarios: subject-driven generation and controllable generation. Our extensive empirical evaluation shows that OFT yields marked advancements over prior methods, delivering higher generation quality, faster convergence rates, and more stable finetuning dynamics. Furthermore, OFT demonstrates increased sample efficiency by requiring considerably fewer images and training epochs to reach effective convergence.
    \item For the controllable generation task, we propose a novel control consistency metric that quantitatively assesses the controllability of the model. The metric is computed by extracting the control signal from a generated image and comparing it to the original control input, thereby providing additional evidence for OFT’s superior controllability.
\end{itemize}

\section{Related Work}

\textbf{Text-to-image diffusion models}. Recent years have witnessed substantial advances in text-to-image synthesis~\cite{saharia2022photorealistic,ramesh2022hierarchical,rombach2022high,gu2022vector,nichol2022glide}, propelled by innovations in diffusion-based generative modeling~\cite{ho2020denoising,song2021score,song2021denoising,dhariwal2021diffusion} alongside developments in multimodal representation learning~\cite{su2020vlbert,lu2019vilbert,radford2021learning,wang2021simvlm,li2022blip,li2023blip,alayrac2022flamingo,schuhmann2022laion}. Architectures such as GLIDE~\cite{nichol2022glide} and Imagen~\cite{saharia2022photorealistic} operate on pixel-space diffusion models. Specifically, GLIDE jointly optimizes a text encoder with a diffusion prior using aligned text-image pairs, while Imagen opts for a pretrained, frozen text encoder. Meanwhile, Stable Diffusion~\cite{rombach2022high} and DALL-E2~\cite{ramesh2022hierarchical} deploy diffusion models within latent spaces: Stable Diffusion relies on VQ-GAN~\cite{esser2021taming} for its latent visual codebook, and DALL-E2 utilizes CLIP~\cite{radford2021learning} to craft a shared embedding space for both images and text. Beyond diffusion models, alternative paradigms such as generative adversarial networks~\cite{reed2016generative,zhang2017stackgan,xu2018attngan,li2019controllable} and autoregressive models~\cite{ramesh2021zero,ding2021cogview,wu2022nuwa,yuscaling} have also addressed the text-to-image task. OFT represents a model-agnostic finetuning mechanism, applicable across diverse text-to-image diffusion architectures.

\textbf{Subject-driven generation}. Several approaches, including \cite{avrahami2022blended,nichol2022glide}, integrate user-provided masks as auxiliary conditions to avoid altering the main subject. Inversion-based strategies~\cite{choi2021ilvr,dhariwal2021diffusion,ramesh2022hierarchical,gal2022image} can edit image backgrounds without impacting the subject. Moreover, \cite{hertz2022prompt} achieves localized and global modifications independent of explicit mask inputs. However, these techniques generally struggle to retain fine identity details of the subject. The Pivotal Tuning technique~\cite{roich2022pivotal} finetunes the generative model around an initially inverted latent code, using an explicit regularization to conserve subject identity. Along similar lines, \cite{nitzan2022mystyle} develops a person-specific face generator from a small set of face images. Instance diversity methods such as \cite{casanova2021instance} can yield different instance variations but may compromise unique instance details. DreamBooth~\cite{ruiz2023dreambooth} uses a few subject-specific images and an associated custom token to finetune text-to-image diffusion models with a combination of reconstruction and class-preserving losses. OFT is implemented within the DreamBooth protocol; however, unlike traditional approaches that apply simple finetuning with a reduced learning rate, OF

\textbf{Controllable generation}. Image-to-image translation can be interpreted as a controllable generative process, for which earlier approaches predominantly rely on conditional generative adversarial networks~\cite{isola2017image,zhu2017toward,wang2018high,park2019semantic,choi2018stargan}. In addition, diffusion models have also been incorporated into image-to-image translation pipelines~\cite{saharia2022palette,voynov2022sketch,wang2022pretraining}. More recent advancements include ControlNet~\cite{zhang2023adding}, which adapts pretrained diffusion models through finetuning, leveraging additional control signals to attain remarkable controllable generation capability. Similarly, the contemporaneous T2I-Adapter~\cite{mou2023t2i} employs finetuning of a pretrained diffusion model to enhance the controllability of synthesized images. In line with the experimental protocols of \cite{zhang2023adding,mou2023t2i}, we introduce OFT as a finetuning mechanism for pretrained diffusion models, which results in consistently improved control using fewer training samples and a reduced number of finetuned parameters. Notably, OFT does not increase computational requirements during inference.

\textbf{Model finetuning}. The practice of finetuning large-scale pretrained models on downstream applications has gained substantial traction in recent years~\cite{devlin2018bert,brown2020language,he2022masked}. Within this context, adaptation techniques (\emph{e.g.}, \cite{hulora2022,houlsby2019parameter,pfeiffer2020adapterfusion}) are extensively explored in natural language processing tasks. Of particular relevance to OFT is LoRA~\cite{hulora2022}, which posits a low-rank update structure added to weights during finetuning. Distinctively, OFT employs a layer-wise shared orthogonal transformation for multiplicative adjustment of neuron weights, which analytically maintains pairwise angular relationships among neurons within a layer, thereby providing enhanced training stability.

\section{Orthogonal Finetuning}

\subsection{Why Does Orthogonal Transformation Make Sense?}

We begin by elucidating the rationale for employing orthogonal transformations in the finetuning of text-to-image diffusion models. The discussion is organized around two fundamental questions: (1) the motivation for adjusting the angular orientation of neurons (\emph{i.e.}, their direction), and (2) the reasons for utilizing orthogonal transformation for such angular finetuning.

Addressing the first question, we are motivated by empirical findings from \cite{liu2018decoupled,chen2020angular}, which indicate that differences in angular features are indicative of semantic disparities. According to SphereNet~\cite{liu2017deep}, constraining all neurons within a neural network to reside on a unit hypersphere maintains representational capacity and often enhances generalization. This suggests that the essential information in data is primarily encoded in the neuron directions. To further underscore the relevance of neuronal angles, we design a simple experiment depicted in Figure~\ref{fig:toy_ae}, where a conventional convolutional autoencoder is trained on a dataset of flower images. During training, the network uses the standard inner product for computing feature maps ($z$ representing the output element from a convolution filter $\bm{w}$ and $\bm{x}$ as the windowed input). For testing, we assess three approaches for constructing the feature map: (a) using the original inner product, (b) retaining only magnitude information, and (c) retaining only angular information. The findings, presented in Figure~\ref{fig:toy_ae}, demonstrate that utilizing only the neurons' angular information almost perfectly reconstructs the input images, whereas magnitude alone is uninformative. It is noteworthy that throughout training, no cosine-based activation (c) was employed; the model was trained exclusively via the inner product. This experiment illustrates that neuronal direction (angle) predominantly preserves semantic content, implying that adjusting neuron directions is likely the most effective strategy for manipulating image semantics.

Concerning the second question, the most direct strategy for finetuning neuronal directions is to rotate or reflect all neurons within a layer simultaneously, leading naturally to the use of orthogonal transformations. Although other angular modifications—such as rotating each neuron by an individual angle—could offer additional flexibility, we find that orthogonal transformations provide a balanced compromise between adaptability and regularization. Supporting this, \cite{liu2021orthogonal} demonstrates that orthogonal transformations offer substantial expressive power for training neural networks. To substantiate our perspective, we conduct an experiment that assesses the regularizing effect induced by enforcing orthogonality. Leveraging the experimental framework of ControlNet~\cite{zhang2023adding}, we present our results in Figure~\ref{fig:ortho}. Here, standard OFT is shown to be both robust and accurate in control after 20 epochs of training. Conversely, removing the orthogonality constraint leads to a complete failure to generate realistic images or achieve meaningful control. This empirical evidence highlights the critical role of orthogonality within OFT.

\subsection{General Framework}

A standard approach to finetuning involves adjusting a model—or selected layers—using gradient descent with a reduced learning rate. This low learning rate effectively restricts updates, promoting only minor changes from the pretrained parameters, which can be formulated as minimizing $\thickmuskip=2mu \medmuskip=2mu \| \bm{M} - \bm{M}^0\|$, where $\bm{M}$ represents the finetuned weights and $\bm{M}^0$ those of the pretrained model. Although this regularization is intuitively reasonable, it might still allow excessive flexibility, especially when adapting large-scale models. To tighten this control, LoRA enforces a further constraint by restricting the rank of the weight modification: $\thickmuskip=2mu \medmuskip=2mu \text{rank}(\bm{M} - \bm{M}^0)=r'$, where $r'$ is a designated small integer. In contrast, OFT imposes a restriction on the neuron-wise relational structure by requiring $\thickmuskip=2mu \medmuskip=2mu \| \text{HE}(\bm{M})-\text{HE}(\bm{M}^0) \|=0$, where $\text{HE}(\cdot)$ captures the hyperspherical energy associated with a weight matrix.

To exemplify this, consider a fully connected layer parameterized as $\thickmuskip=2mu \medmuskip=2mu \bm{W}=\{\bm{w}_1,\cdots,\bm{w}_n\}\in\mathbb{R}^{d\times n}$, where each column $\bm{w}_i\in\mathbb{R}^d$ corresponds to the $i$-th neuron's weights, and $\bm{W}^0$ is the pretrained analogue. For input $\thickmuskip=2mu \medmuskip=2mu \bm{x}\in\mathbb{R}^d$, the resulting output is given by $\thickmuskip=2mu \medmuskip=2mu \bm{z}=\bm{W}^\top\bm{x}$. OFT may thus be interpreted as minimizing discrepancies in hyperspherical energy between updated and pretrained models:

\begin{equation}\label{eq:oft_principle}
\footnotesize
\min_{\bm{W}} \left\| \text{HE}(\bm{W})-\text{HE}(\bm{W}^0)\right\|~~\Leftrightarrow~~\min_{\bm{W}} \bigg{\|} \sum_{i\neq j} \|\hat{\bm{w}}_i-\hat{\bm{w}}_j\|^{-1}-\sum_{i\neq j} \|\hat{\bm{w}}^0_i-\hat{\bm{w}}^0_j\|^{-1}\bigg{\|}
\end{equation}

where each normalized neuron vector is $\thickmuskip=2mu \medmuskip=2mu \hat{\bm{w}}_i:=\bm{w}_i/\|\bm{w}_i\|$, and the hyperspherical energy for a given layer $\bm{W}$ is defined as $\thickmuskip=2mu \medmuskip=2mu \text{HE}(\bm{W}):= \sum_{i\neq j} \|\hat{\bm{w}}_i-\hat{\bm{w}}_j\|^{-1}$. Notably, the expression in Eq.~\eqref{eq:oft_principle} achieves its global minimum of zero when $\bm{W}$ differs from $\bm{W}^0$ solely by a rotation or reflection, i.e., when $\thickmuskip=2mu \medmuskip=2mu \bm{W}=\bm{R}\bm{W}^0$ for some orthogonal matrix $\thickmuskip=2mu \medmuskip=2mu \bm{R}\in\mathbb{R}^{d\times d}$. (Here, $\det(\bm{R})=1$ denotes rotation and $\det(\bm{R})=-1$ corresponds to reflection.) This embodies the core concept behind OFT: it suffices to optimize over orthogonal transformations, learning a shared orthogonal matrix per layer to finetune the network. More precisely, given a pretrained layer $\thickmuskip=2mu \medmuskip=2mu\bm{W}^0\in\mathbb

\begin{equation}\label{eq:oft_general}
    \footnotesize
    \bm{z}=\bm{W}^\top\bm{x}=(\bm{R}\cdot\bm{W}^0)^\top\bm{x},~~~\text{s.t.}~\bm{R}^\top\bm{R}=\bm{R}\bm{R}^\top=\bm{I}
\end{equation}

In this expression, $\bm{W}$ corresponds to the weight matrix updated through OFT fine-tuning, while $\bm{I}$ is the identity matrix. A schematic overview of OFT can be found in Figure~\ref{fig:block}. As with the zero initialization utilized in LoRA, it is crucial for OFT to start fine-tuning from the initial pretrained parameters $\bm{W}^0$. To satisfy this requirement, we set the initial value of the orthogonal matrix $\bm{R}$ to the identity matrix, ensuring the initial model matches the pretrained weights. To maintain the orthogonality of $\bm{R}$ during training, differential orthogonalization methods described in \cite{liu2021orthogonal,lezcano2019cheap} can be leveraged. We elaborate on approaches for maintaining orthogonality that are computationally efficient.

\subsection{Efficient Orthogonal Parameterization}\label{sect:eff_ortho}

Conventional approaches for enforcing orthogonality, such as the differentiable Gram-Schmidt process, are often computationally burdensome and impractical for large models~\cite{liu2021orthogonal}. To mitigate this issue, we employ the Cayley parameterization to construct the orthogonal matrix. Specifically, the matrix $\bm{R}$ is obtained using the formulation $\thickmuskip=2mu \medmuskip=2mu \bm{R}=(\bm{I}+\bm{Q})(I-\bm{Q})^{-1}$, where $\bm{Q}$ is a skew-symmetric matrix, i.e., $\thickmuskip=2mu \medmuskip=2mu \bm{Q}=-\bm{Q}^\top$. The Cayley transform offers substantial computational benefits but is limited in that it generates orthogonal matrices with determinant $1$, i.e., members of the special orthogonal group. Nonetheless, empirical evidence suggests this restriction does not negatively impact performance. Even when using the Cayley transform, parameterizing a full-size orthogonal matrix $\bm{R}$ becomes parameter-inefficient for large $d$. To reduce this cost, we introduce a block-diagonal parameterization for $\bm{R}$, partitioning it into $r$ independent blocks. This results in the following structure:

\begin{equation}
    \footnotesize
    \bm{R}=\text{diag}(\bm{R}_1,\bm{R}_2,\cdots,\bm{R}_r)=\begin{bmatrix}
    \bm{R}_1\in \text{O}(\frac{d}{r}) &  &\\
    &\ddots & \\
    & & \bm{R}_r\in \text{O}(\frac{d}{r})
    \end{bmatrix}\in \text{O}(d)
\end{equation}

Here, $\text{O}(d)$ represents the orthogonal group in $d$ dimensions, with $\thickmuskip=2mu \medmuskip=2mu \bm{R}\in\mathbb{R}^{d\times d}$ and, for each $i$, $\thickmuskip=2mu \medmuskip=2mu \bm{R}_i\in\mathbb{R}^{d/r\times d/r}$ denoting orthogonal matrices. In the specific case where $\thickmuskip=2mu \medmuskip=2mu r=1$, the block-diagonal orthogonal matrix simply reduces to an unconstrained orthogonal matrix. Considering an orthogonal matrix of dimensions $\thickmuskip=2mu \medmuskip=2mu d\times d$, there are $\thickmuskip=2mu \medmuskip=2mu d(d-1)/2$ free parameters, so the computational complexity is $\mathcal{O}(d^2)$. By contrast, an $r$-block diagonal orthogonal matrix comprises $\thickmuskip=2mu \medmuskip=2mu d(d/r-1)/2$ parameters, and thus its complexity is reduced to $\thickmuskip=2mu \medmuskip=2mu \mathcal{O}(d^2/r)$. To decrease parameters even further, one can enforce parameter sharing between blocks, i.e., let $\thickmuskip=2mu \medmuskip=2mu\bm{R}_i=\bm{R}_j$ for all $i\neq j$. In this setting, the parameter complexity is further lowered to $\mathcal{O}(d^2/r^2)$. Regardless of these reductions in parameterization, the overall matrix $\bm{R}$ always maintains orthogonality, thereby preserving hyperspherical energy without compromise.

We now compare the parameter efficiency of OFT and LoRA. For LoRA with a low-rank dimension $r'$, the number of trainable parameters is $\thickmuskip=2mu \medmuskip=2mu r'(d+n)$. When $r$ and $r'$ are treated as dependent on $d$ (for example, $\thickmuskip=2mu \medmuskip=2mu r=r'=\alpha d$ for some fixed $0<\alpha\leq 1$), LoRA has a parameter complexity of $\mathcal{O}(d^2+dn)$, whereas OFT attains $\mathcal{O}(d)$. To clarify these differences, consider the following scenario: for a weight matrix of size $\thickmuskip=2mu \medmuskip=2mu 128\times 128$, setting $\thickmuskip=2mu \medmuskip=2mu r'=8$ results in LoRA having $2,048$ parameters, while with $\thickmuskip=2mu \medmuskip=2mu r=8$ and no shared blocks, OFT contains $960$ parameters.

\subsection{Constrained Orthogonal Finetuning}

To restrict the degree of freedom in OFT, we can require that the updated model stays within a restricted vicinity of the pretrained parameters. Specifically, COFT employs the following forward propagation formula:

\begin{equation}\label{eq:coft_general}
    \footnotesize
    \bm{z}=\bm{W}^\top\bm{x}=(\bm{R}\cdot\bm{W}^0)^\top\bm{x},~~~\text{s.t.}~\bm{R}^\top\bm{R}=\bm{R}\bm{R}^\top=\bm{I},~\norm{\bm{R}-\bm{I}}\leq \epsilon
\end{equation}

This formulation enforces both the orthogonality constraint and an $\epsilon$-bounded deviation from the identity matrix. Orthogonality is enforced via the Cayley parameterization discussed earlier in Section~\ref{sect:eff_ortho}. Nevertheless, adding the $\epsilon$-deviation restriction to the Cayley-parameterized matrix is a nontrivial task. To better understand the effect of the Cayley transform, we employ the Neumann series to approximate $\thickmuskip=2mu \medmuskip=2mu \bm{R}=(\bm{I}+\bm{Q})(I-\bm{Q})^{-1}$, giving $\thickmuskip=2mu \medmuskip=2mu \bm{R}\approx

series converges with respect to the operator norm). As a result, the constraint $\thickmuskip=2mu \medmuskip=2mu\|\bm{R}-\bm{I}\|\leq\epsilon$ can be internalized within the Cayley transform, leading to an analogous condition $\thickmuskip=2mu \medmuskip=2mu \|\bm{Q}-\bm{0}\|\leq\epsilon'$, where $\bm{0}$ is the zero matrix and $\epsilon'$ is a new tuning parameter, distinct from $\epsilon$. Enforcing this condition on $\bm{Q}$ is straightforward using projected gradient descent techniques. To ensure that the orthogonal matrix $\bm{R}$ is initialized as the identity, we initialize $\bm{Q}$ to all zeros. COFT thus effectively imposes two distinct constraints: one on hyperspherical energy difference and another limiting divergence from the pretrained weights. The latter is generally used implicitly in other finetuning approaches, but COFT makes it explicit. We find that, even though OFT performs strongly, the explicit regularization of deviation in COFT leads to more stable finetuning outcomes compared to OFT. Figure~\ref{fig:eps_coft} shows how the parameter $\epsilon$ influences COFT performance. Specifically, as $\epsilon$ increases, COFT's behavior converges to that of OFT; with decreasing $\epsilon$, the behavior aligns more closely with the pretrained text-to-image diffusion network.

\subsection{Re-scaled Orthogonal Finetuning}

We introduce a straightforward modification to standard OFT, allowing each neuron to learn an additional scaling factor. The rationale is that rescaling neuron outputs leaves their hyperspherical energy unaffected, since magnitude is normalized. The forward computation now becomes $\thickmuskip=2mu \medmuskip=2mu\bm{z}=(\bm{R}\bm{W}^0\bm{D})^\top\bm{x}$\footnote{\textbf{Errata}: The NeurIPS camera-ready paper incorrectly wrote the forward pass for re-scaled OFT as $\thickmuskip=2mu \medmuskip=2mu\bm{z}=(\bm{D}\bm{R}\bm{W}^0)^\top\bm{x}$. Our original implementation used the correct form, so the results in Appendix~\ref{app:rescaled} remain valid.} where $\thickmuskip=2mu \medmuskip=2mu \bm{D}=\text{diag}(s_1,\cdots,s_n)\in\mathbb{R}^{n\times n}$, with all diagonal entries $s_1,\ldots,s_n$ as positive, trainable parameters. Unlike the OFT formulation in Eq.~\eqref{eq:oft_general}, where only $\bm{R}$ is updated, the re-scaled approach makes both $\bm{D}$ and $\bm{R}$ learnable. This re-scaling increases OFT's representational power, yet introduces only a minimal parameter overhead. For our main experiments, we use the original OFT to isolate the benefits of orthogonal transformations. Nonetheless, empirical observations indicate that re-scaled OFT typically provides superior results (refer to Appendix~\ref{app:rescaled}).

\section{Intriguing Insights and Discussions}

\textbf{OFT's architectural independence}. OFT is, by design, compatible with virtually any neural architecture. In the case of Transformers, LoRA is conventionally applied to the attention parameters~\cite{hulora2022}. To ensure a fair baseline for comparison with LoRA, our application of OFT is likewise restricted to the attention weights in experiments. Beyond dense layers, OFT is also readily applicable to convolutional layers. The block-diagonal nature of $\bm{R}$ reveals meaningful interpretations when used with convolutions, a feature that differentiates it from LoRA. Assigning each block to correspond to an input channel, the blocks operate on unique channels, mirroring depth-wise convolution kernel learning~\cite{chollet2017xception}. Alternatively, if the same orthogonal block is applied across all channels, the neurons are jointly transformed using a shared orthogonal matrix. We provide an initial investigation of OFT applied to convolutional layer finetuning in Appendix~\ref{app:convolution}.

\textbf{Connection to LoRA}. LoRA incorporates a low-rank matrix to ensure that the pretrained weight matrix's information is largely preserved and does not undergo drastic changes. In contrast, OFT imposes an orthogonality constraint—maintaining full rank—on the transformation applied to the pretrained weights, thus safeguarding the pretraining content. The OFT forward pass can be represented as $\thickmuskip=2mu \medmuskip=2mu \bm{z}=(\bm{R}\bm{W}^0)^\top\bm{x}=(\bm{W}^0+(\bm{R}-\bm{I})\bm{W}^0)^\top\bm{x}$, where the term $\thickmuskip=2mu \medmuskip=2mu (\bm{R}-\bm{I})\bm{W}^0$ mirrors the role of LoRA's low-rank update. Since $\bm{W}^0$ is generally full rank, OFT will also perform a low-rank update whenever $\thickmuskip=2mu \medmuskip=2mu \bm{R}-\bm{I}$ is of low rank. Analogous to LoRA’s rank parameter $r'$, OFT utilizes a block diagonal parameter $r$ to constrain the quantity of parameters that must be learned. Importantly, LoRA and OFT achieve parameter efficiency via fundamentally different principles. LoRA leverages low-rank approximations, whereas OFT reduces trainable parameters by exploiting a sparse block-diagonal orthogonal structure.

\textbf{Why OFT converges faster}. Firstly, as illustrated in Figure~\ref{fig:toy_ae}, the principal modification for altering the representational semantics involves adjusting neuron directions—a mechanism that OFT is specifically constructed to perform. Moreover, OFT can be understood as adapting neurons on a smooth hyperspherical manifold, a property that leads to more favorable optimization dynamics. This observation is also supported empirically by the findings in \cite{liu2021orthogonal}.

\textbf{Why not minimize hyperspherical energy}. Distinct from the methodology in \cite{liu2021orthogonal}, our approach does not pursue minimization of hyperspherical energy. In the context of classification, the objective is to avoid redundant neurons. Achieving minimum hyperspherical energy would enforce that neurons are uniformly distributed over the hypersphere, which does not benefit the finetuning process and may actually compromise the pretrained representations.

\textbf{Balancing flexibility and regularity in finetuning}. Our findings highlight a fundamental tension between flexibility and regularity. Conventional finetuning provides the highest degree of flexibility but suffers from limited stability and a heightened risk of model collapse. OFT, despite its conceptual simplicity, achieves an effective compromise between these two objectives by maintaining the pairwise angles between neurons. The block-diagonal structure serves as a stricter form of regularization on the orthogonal matrix.

\textbf{Inference efficiency maintained}. In contrast to ControlNet, the OFT approach imposes no additional computational cost during inference on the adapted model. At inference time, the learned orthogonal matrix $\bm{R}$ can be merged with the pretrained weights $\bm{W}^0$ through matrix multiplication, resulting in the updated weights $\thickmuskip=2mu \medmuskip=2mu \bm{W}=\bm{R}\bm{W}^0$. This operation ensures the inference speed remains unchanged from that of the original pretrained model.

\section{Experiments and Results}

\textbf{Experimental setup}. All experiments leverage Stable Diffusion v1.5~\cite{rombach2022high} as the foundation for text-to-image generation. To maintain impartiality, samples from all compared methods are selected at random. The procedure for subject-driven generation is largely modeled after DreamBooth~\cite{ruiz2023dreambooth}, whereas protocols for controllable generation draw from ControlNet~\cite{zhang2023adding} and T2I-Adapter~\cite{mou2023t2i}. For a fair comparison with LoRA, both OFT and COFT are restricted to being applied at the same layers as LoRA. Additional methodological details and experimental results are included in Appendix~\ref{app:settings}.

\subsection{Subject-driven Generation}

\textbf{Configuration}. Baseline methods include DreamBooth~\cite{ruiz2023dreambooth} and LoRA~\cite{hulora2022}. All approaches employ the loss function as specified in DreamBooth. For DreamBooth and LoRA, the experimental protocols closely follow the respective original works, employing optimal hyperparameter configurations as recommended. Supplemental results and discussion are provided in Appendix~\ref{app:settings},\ref{app:coft_oft},\ref{app:more_qualitative},\ref{app:failure}.

\textbf{Finetuning stability and convergence}. To begin, we assess the stability and convergence rate during finetuning for DreamBooth, LoRA, OFT, and COFT. Figures~\ref{fig:teaser} and~\ref{fig:db_convergence} present the corresponding results. Notably, both OFT and COFT maintain stable adaptation of the diffusion model throughout the process. By 400 iterations, DreamBooth and the OFT-based approaches are able to provide effective guidance, while LoRA is unable to reliably retain the identity of the subject. At the 2000-iteration mark, DreamBooth starts to suffer from image collapse, and LoRA is unable to render a yellow shirt correctly, often substituting it with yellow fur. In contrast, OFT and COFT continue to demonstrate robust and steady control over the generated outputs. These outcomes demonstrate that OFT offers rapid as well as reliable convergence in subject-driven generation tasks. It is worth emphasizing that robustness with respect to finetuning iteration count is a significant advantage, as it reduces the need for extensive hyperparameter tuning across subjects. Both OFT and COFT permit the use of larger iteration counts directly, without necessitating sensitive adjustment. Particularly for COFT, with an appropriate choice of $\epsilon$, setting the learning rate and iteration number becomes a straightforward process.

\textbf{Quantitative comparison}. Following the methodology of \cite{ruiz2023dreambooth}, we perform a quantitative analysis in terms of subject fidelity (DINO~\cite{caron2021emerging}, CLIP-I~\cite{radford2021learning}), alignment with text prompts (CLIP-T~\cite{radford2021learning}), and diversity of samples (LPIPS~\cite{zhang2018unreasonable}). The CLIP-I metric computes the mean cosine similarity across CLIP embedding pairs from generated and actual subject images. DINO works analogously to CLIP-I, substituting in ViT S/16 DINO representations. For CLIP-T, the metric measures the average cosine similarity between the text prompt and generated images using CLIP embeddings. Additionally, average LPIPS cosine similarity quantifies diversity across generated images with the same prompt for a given subject. The data in Table~\ref{table:dreambooth_final} indicate that both COFT and OFT surpass DreamBooth and LoRA substantially on the DINO and CLIP-I benchmarks, while also delivering marginally superior or comparable outcomes in prompt consistency and diversity scores. To address randomness, each method is finetuned 30 times using different seeds but the same pretrained weights. Collectively, the findings establish that our OFT strategy achieves not only stronger convergence characteristics and stability, but also superior and consistently repeatable overall performance.

\textbf{Qualitative comparison}. To provide clearer insights into the advantages of OFT, we present several randomly selected samples of subject-driven image generation in Figure~\ref{fig:dreambooth_final}. For consistency and fairness, all outputs are produced from the same finetuned model with each respective method, and no prompts are specifically adjusted or optimized for individual images. For each technique, the model version that attains the highest validation CLIP score is chosen for illustration. As displayed in Figure~\ref{fig:dreambooth_final}, both OFT and COFT are able to robustly maintain the semantic integrity of the subject, in contrast to LoRA, which frequently fails to preserve subject identity (\emph{e.g.}, in the bowl example, LoRA entirely fails to retain the correct subject). Furthermore, OFT and COFT demonstrate significantly enhanced accuracy in following text prompts, whereas DreamBooth, despite often capturing subject identity well, struggles to ensure generated images conform to the textual guidance (\emph{e.g.}, the bowl case in the first row). These qualitative findings indicate that our OFT approach offers superior balance between control over image attributes and the faithful representation of subject identity. Additionally, OFT's performance remains stable even when the number of iterations is high, in contrast to both DreamBooth and LoRA, which degrade under similar conditions. Further qualitative examples can be found in Appendix~\ref{app:more_qualitative}. A human study in Appendix~\ref{app:human} further substantiates OFT's performance advantage.

\subsection{Controllable Generation}

\textbf{Settings}. For baselines, we include ControlNet~\cite{zhang2023adding}, T2I-Adapter~\cite{mou2023t2i}, and LoRA~\cite{hulora2022}. We address three notably challenging controllable generation benchmarks in this work: canny edge to image~(C2I) conversion using the COCO dataset~\cite{lin2014microsoft}, segmentation map to image~(S2I) on ADE20K~\cite{zhou2017scene}, and landmark to face~(L2F) synthesis using the CelebA-HQ dataset~\cite{karras2018progressive,xia2021tedigan}. Each baseline method, as well as ours, is used to finetune Stable Diffusion~(SD) v1.5 on the respective dataset for 20 epochs. Additional results are provided in Appendix~\ref{app:more_qualitative},\ref{app:more_control_task},\ref{app:failure}.

\textbf{Convergence}. To assess the convergence behavior, we compare ControlNet, T2I-Adapter, LoRA, and COFT within the context of the L2F task. Both quantitative and qualitative analyses are provided. For the quantitative metric, we calculate the mean $\ell_2$ distance between the control face landmarks and those predicted by the model. Figure~\ref{fig:face_convergence} illustrates how the face landmark error evolves as each model is finetuned over varying epoch counts. The results indicate that COFT and OFT reach convergence substantially faster than their counterparts. Notably, LoRA requires 20 epochs to reach the performance level that OFT achieves by epoch 8. Importantly, both OFT and COFT utilize a comparable number of trainable parameters as LoRA—which is substantially less than that of ControlNet—but demonstrate markedly superior convergence speed compared to existing approaches. Moreover, the accelerated convergence exhibited by OFT is further substantiated by the outcomes depicted in Figure~\ref{fig:teaser}. In particular, the rightmost example in Figure~\ref{fig:teaser} demonstrates OFT's superior data efficiency relative to ControlNet and LoRA, as it achieves effective convergence using just 5\% of the ADE20K dataset. For the qualitative comparison, emphasis is placed on OFT, COFT, and ControlNet, given that ControlNet yields the closest landmark error to OFT. As depicted in Figure~\ref{fig:face_convergence_qual}, both OFT and COFT exhibit steady convergence, with generated face poses incrementally aligning with the prescribed control landmarks. In contrast, ControlNet displays an abrupt emergence of controllability following the 8th epoch—echoing the sharp convergence transition reported in \cite{zhang2023adding}. The smoother and more consistent convergence trajectory of OFT facilitates practical use, as its training process is both predictable and stable, thereby simplifying the tuning of hyperparameters.

\textbf{Quantitative comparison}. To assess the effectiveness of controllable generation, we propose a control consistency metric. This metric operates by extracting the control signal from the generated output and measuring its alignment with the original input control. For the C2I task, we utilize IoU and F1 score, whereas for the S2I task, we report mean IoU along with mean and overall accuracy. For the L2F task, the mean $\ell_2$ distance between the original control landmarks and the corresponding predicted landmarks is used. Additional specifics regarding these consistency measurements can be found in Appendix~\ref{app:settings}. In our experiments, each method is evaluated using its optimal hyperparameter configuration. According to the results in Table~\ref{table:control_gen}, both OFT and COFT significantly surpass other techniques in terms of control precision and strength. Adapter-based models such as T2I-Adapter and ControlNet tend to learn more slowly and attain lower final performance. LoRA exceeds ControlNet on the S2I task but trails behind in both C2I and L2F tasks. Overall, our findings suggest that LoRA’s maximum attainable performance remains limited, even under rigorous hyperparameter optimization. In contrast, OFT continues to exhibit improvements as we increase the number of trainable parameters, indicating its performance has not plateaued. Notably, our quantitative methodology for evaluating controllable generation stands as a significant innovation, as it provides precise assessment of control capabilities in fine-tuned models, constrained only by the performance of the employed segmentation or detection backbone.

\textbf{Qualitative comparison}. We further perform a qualitative assessment of OFT and COFT against leading existing approaches such as ControlNet, T2I-Adapter, and LoRA. As presented in Figure~\ref{fig:control_final}, randomly sampled outputs indicate that both OFT and COFT deliver visually convincing and high-quality images while maintaining precise controllability. In the S2I setting, it is evident that LoRA fails to conform to the given segmentation maps, whereas ControlNet, OFT, and COFT faithfully adhere to the input structures. Notably, when compared with ControlNet, OFT and COFT not only preserve a high degree of fidelity and intricate detail, but also achieve more coherent and plausible geometric composition, despite having a significantly smaller model size. With respect to the C2I task, OFT and COFT are able to synthesize semantically relevant images using coarse Canny edge inputs, surpassing the outputs of T2I-Adapter and LoRA, which exhibit inferior performance. For the L2F task, our methods provide the most reliable control over facial pose generation, even under difficult pose conditions. Across all three categories of controllable generation, OFT and COFT consistently outperform state-of-the-art baselines in qualitative terms, validating the efficacy of our OFT framework for controllable image synthesis. To offer a broader perspective, we present additional qualitative results for these control tasks in Appendix~\ref{app:control_exp}, and further demonstrate the applicability of OFT to various other control scenarios—including dense pose to body synthesis, sketch to image translation, and depth-to-image generation—in Appendix~\ref{app:more_control_task}.

\section{Concluding Remarks and Open Problems}

Drawing on the insight that angular relationships among neurons play a critical role in shaping visual semantics, we introduce a straightforward yet powerful finetuning strategy—orthogonal finetuning—to enhance control in text-to-image diffusion models. In particular, our approach addresses two primary tasks: subject-driven generation and controllable generation. When evaluated against existing techniques, OFT achieves greater controllability and more robust finetuning stability, all while requiring a smaller set of finetuning parameters. A key advantage of OFT is that it does not increase inference costs, resulting in a solution that is both efficient and practical for deployment.

However, OFT gives rise to several intriguing open questions. Firstly, while the Cayley parametrization employed by OFT ensures orthogonality, it relies on matrix inversion, posing constraints on the method's scalability. We partially mitigate this by adopting a block diagonal parametrization, yet devising a differentiable, accelerated approach for matrix inversion is an unresolved challenge. Secondly, OFT inherently supports compositionality; specifically, the orthogonal matrices generated by various OFT finetuning procedures can be combined through multiplication, yielding another orthogonal matrix. An open area of inquiry is whether this composition retains the full scope of task-specific knowledge across all downstream applications. Finally, the model's parameter efficiency is primarily linked to the block diagonal design, which can introduce structural bias and constrain adaptability. Discovering strategies to enhance parameter efficiency that reduce bias and preserve flexibility remains an important research direction.